\clearpage
\begin{explanationbox}
  \textbf{\textcolor{CExplanation}{一、为什么要单独研究幂函数？}}
\begin{center}\begin{tikzpicture}[font=\small,>=stealth,line join=round]
\draw[gray!40,rounded corners=3pt] (0,0) rectangle (14.2,8.3);\node at (7.1,7.700000000000001) {同过(1,1)，底数跨过1时，上下顺序互换};
\begin{scope}[shift={(1.6625,2.068181818181818)},x=1.8750000000000002cm,y=0.8363636363636363cm]\draw[->,gray] (-0.3,0)--(2.9,0) node[right]{$x$};\draw[->,gray] (0,-0.5)--(0,5) node[above]{$y$};\draw[blue,thick,solid,domain=0.03:2.15,samples=180] plot(\x,{(\x)^2});\draw[teal,thick,dashed,domain=0.03:2.6,samples=180] plot(\x,{\x});\draw[violet,thick,densely dotted,domain=0.03:2.6,samples=180] plot(\x,{sqrt(\x)});\fill[orange!85!black] (1,1) circle(2.2pt);\draw[gray,dashed] (1,0)--(1,4.6);\node[below=5pt] at (1,0){$1$};\end{scope}\node[anchor=west,align=left,text width=5.0cm] at (8.4,4.0) {蓝实线：$x^2$\\[8pt]青虚线：$x$\\[8pt]紫点线：$\sqrt{x}$\\[8pt]$0<x<1$：$x^2<x<\sqrt{x}$\\[8pt]$x>1$：$x^2>x>\sqrt{x}$};
\node[text width=12.6cm,align=center] at (7.1,.65) {这里始终固定系数为1；三条曲线在正半轴上只有公共点(1,1)。};
\end{tikzpicture}\end{center}



  初中已经学过正比例函数 $y=x$、二次函数 $y=x^2$、反比例函数 $y=1/x$；它们看起来很不同，却可以统一写成 $y=x^\alpha$，指数分别是 $1,2,-1$。再加入 $y=\sqrt{x}=x^{1/2}$，就能发现：只改变一个指数，便能得到很多不同的图像。

  \medskip
  
\end{explanationbox}
\clearpage
\begin{explanationbox}
\textbf{\textcolor{CExplanation}{二、为什么幂函数前面的系数必须是 $1$？}}
\begin{center}\begin{tikzpicture}[font=\small,>=stealth,line join=round]
\draw[gray!40,rounded corners=3pt] (0,0) rectangle (14.2,8.3);\node at (7.1,7.700000000000001) {固定指数2，改变系数会改变每一点的高度};
\begin{scope}[shift={(4.1,2.11)},x=1.6666666666666665cm,y=0.9199999999999999cm]\draw[->,gray] (-1.8,0)--(1.8,0) node[right]{$x$};\draw[->,gray] (0,-0.5)--(0,4.5) node[above]{$y$};\draw[blue,thick,solid,domain=-1.35:1.35,samples=180] plot(\x,{(\x)^2});\draw[teal,thick,dashed,domain=-1.35:1.35,samples=180] plot(\x,{2*(\x)^2});\draw[violet,thick,densely dotted,domain=-1.35:1.35,samples=180] plot(\x,{.5*(\x)^2});\fill[orange!85!black] (1,1) circle(2.2pt);\fill[orange!85!black] (1,2) circle(2.2pt);\fill[orange!85!black] (1,0.5) circle(2.2pt);\draw[gray,dashed] (1,0)--(1,3.9);\node[below] at (1,0){$1$};\end{scope}\node[anchor=west,align=left,text width=5.0cm] at (8.4,4.0) {蓝：$x^2$\\[9pt]青：$2x^2$\\[9pt]紫：$\frac12x^2$\\[9pt]$x=1$处高度为$1,2,\frac12$\\[9pt]严格定义只取系数1};
\node[text width=12.6cm,align=center] at (7.1,.65) {同一指数也能有不同高度；固定系数1，才是单独研究指数的基本幂函数。};
\end{tikzpicture}\end{center}



  这不是由某条定理强迫得到的，而是高中数学对幂函数的定义。这样规定的好处，是把指数对图像的影响与前面系数对图像的影响分开。

  图中把 $x^2$ 的高度乘以 $2$ 或 $1/2$，横坐标始终不变。若乘负数，还要关于 $x$ 轴翻转，见易错九。先研究系数为 $1$ 的基本图像，更容易分清指数与系数的不同作用。

  为什么偏偏选 $1$？因为乘以 $1$ 什么也不改变。对一般式 $y=Ax^\alpha$，$x=1$ 时 $y=A$；固定 $A=1$ 后，所有基本幂函数都从同一个比较点 $(1,1)$ 出发。$2x^2$ 仍然是合法的函数，只是不叫高中严格定义下的幂函数。

\end{explanationbox}
\clearpage
\begin{explanationbox}
\textbf{\textcolor{CExplanation}{三、为什么指数正负决定上升或下降？}}
\begin{center}\begin{tikzpicture}[font=\small,>=stealth,line join=round]
\draw[gray!40,rounded corners=3pt] (0,0) rectangle (14.2,8.8);\node at (7.1,8.200000000000001) {固定指数：正指数保序，负指数反序};
\node at (3.1500000000000004,7.3) {指数固定为2};\begin{scope}[shift={(1.3533333333333335,3.3000000000000003)},x=1.6333333333333335cm,y=0.6cm]\draw[->,gray] (-0.4,0)--(2.6,0) node[right]{$x$};\draw[->,gray] (0,-0.5)--(0,5) node[above]{$y$};\draw[blue,thick,solid,domain=0.1:2.15,samples=180] plot(\x,{(\x)^2});\fill[orange!85!black] (1,1) circle(2pt);\fill[orange!85!black] (2,4) circle(2pt);\node[below=4pt] at (1,0){$1$};\node[below=4pt] at (2,0){$2$};\end{scope}\node[align=center,text width=5.8cm] at (3.1500000000000004,1.7) {$1<2$，而$1^2<2^2$\\底数增大，幂值增大};\node at (9.95,7.3) {指数固定为$-2$};\begin{scope}[shift={(8.153333333333334,3.3000000000000003)},x=1.6333333333333335cm,y=0.6cm]\draw[->,gray] (-0.4,0)--(2.6,0) node[right]{$x$};\draw[->,gray] (0,-0.5)--(0,5) node[above]{$y$};\draw[orange!85!black,thick,solid,domain=0.46:2.3,samples=180] plot(\x,{1/((\x)^2)});\fill[orange!85!black] (1,1) circle(2pt);\fill[orange!85!black] (2,0.25) circle(2pt);\node[below=4pt] at (1,0){$1$};\node[below=4pt] at (2,0){$2$};\end{scope}\node[align=center,text width=5.8cm] at (9.95,1.7) {$1<2$，而$1^{-2}>2^{-2}$\\负指数是倒数，方向反转};
\node[text width=12.6cm,align=center] at (7.1,.65) {同尺度，横轴为底数$x$；图取代表指数，证明适用于任意实数指数。};
\end{tikzpicture}\end{center}



  以正指数 $2$ 为例，$x$ 由 $1,2,4$ 依次变大，$x^2$ 变成 $1,4,16$；而 $x^{-2}=1/x^2$ 依次变成 $1,1/4,1/16$。正指数让底数变大时幂值变大；负指数相当于先取正指数的幂，再取倒数，于是大小关系反过来。零指数则始终等于 $1$。

  \medskip
  
\end{explanationbox}
\clearpage
\begin{explanationbox}
\textbf{\textcolor{CExplanation}{四、为什么同样是上升，曲线却有不同的形状？}}
\begin{center}\begin{tikzpicture}[font=\small,>=stealth,line join=round]
\draw[gray!40,rounded corners=3pt] (0,0) rectangle (14.2,8.5);\node at (7.1,7.9) {正指数三类：都上升，但陡峭程度不同};
\node at (2.4,6.9) {$\alpha=2$};\begin{scope}[shift={(1.0642857142857143,2.9272727272727272)},x=1.2142857142857144cm,y=0.6545454545454545cm]\draw[->,gray] (-0.3,0)--(2.5,0) node[right]{$x$};\draw[->,gray] (0,-0.5)--(0,5) node[above]{$y$};\draw[blue,thick,solid,domain=0.04:2.1,samples=160] plot(\x,{(\x)^2});\fill[orange!85!black] (1,1) circle(2pt);\end{scope}\node[align=center,text width=4cm] at (2.4,1.7) {$y=x^2$\\上升，越来越陡};\node at (7.050000000000001,6.9) {$\alpha=1$};\begin{scope}[shift={(5.714285714285715,2.9272727272727272)},x=1.2142857142857144cm,y=0.6545454545454545cm]\draw[->,gray] (-0.3,0)--(2.5,0) node[right]{$x$};\draw[->,gray] (0,-0.5)--(0,5) node[above]{$y$};\draw[teal,thick,solid,domain=0.04:2.1,samples=160] plot(\x,{\x});\fill[orange!85!black] (1,1) circle(2pt);\end{scope}\node[align=center,text width=4cm] at (7.050000000000001,1.7) {$y=x$\\斜率始终为1};\node at (11.7,6.9) {$\alpha=\frac12$};\begin{scope}[shift={(10.364285714285714,2.9272727272727272)},x=1.2142857142857144cm,y=0.6545454545454545cm]\draw[->,gray] (-0.3,0)--(2.5,0) node[right]{$x$};\draw[->,gray] (0,-0.5)--(0,5) node[above]{$y$};\draw[violet,thick,solid,domain=0.04:2.1,samples=160] plot(\x,{sqrt(\x)});\fill[orange!85!black] (1,1) circle(2pt);\end{scope}\node[align=center,text width=4cm] at (11.7,1.7) {$y=\sqrt{x}$\\上升，越来越平};
\node[text width=12.6cm,align=center] at (7.1,.65) {三图同尺度，只画正半轴；指数1将两种弯曲走势分开。};
\end{tikzpicture}\end{center}



  每次将 $x$ 增加同样的长度，观察 $y$ 改变了多少：
  \begin{center}
  \renewcommand{\arraystretch}{1.25}
  \begin{tabular}{c|cccc}
    $x$ & $1$ & $2$ & $3$ & $4$\\ \hline
    $x^2$ & $1$ & $4$ & $9$ & $16$\\
    $x$ & $1$ & $2$ & $3$ & $4$\\
    $\sqrt{x}$ & $1$ & $\sqrt2$ & $\sqrt3$ & $2$\\
    $1/x$ & $1$ & $1/2$ & $1/3$ & $1/4$
  \end{tabular}
  \end{center}
  $x^2$ 每次增加 $3,5,7$，越来越快；$x$ 每次增加 $1$；$\sqrt{x}$ 每次增加量约为 $0.414,0.318,0.268$，越来越慢；$1/x$ 每次减少量为 $1/2,1/6,1/12$，下降也越来越慢。这几组数帮助我们看懂曲线的形状，完整理由见证明部分。

  \medskip
  
\end{explanationbox}
\clearpage
\begin{explanationbox}
\textbf{\textcolor{CExplanation}{五、为什么在 $1$ 的两边，指数比较方向恰好相反？}}
\begin{center}\begin{tikzpicture}[font=\small,>=stealth,line join=round]
\draw[gray!40,rounded corners=3pt] (0,0) rectangle (14.2,8.3);\node at (7.1,7.700000000000001) {同过(1,1)，底数跨过1时，上下顺序互换};
\begin{scope}[shift={(1.6625,2.068181818181818)},x=1.8750000000000002cm,y=0.8363636363636363cm]\draw[->,gray] (-0.3,0)--(2.9,0) node[right]{$x$};\draw[->,gray] (0,-0.5)--(0,5) node[above]{$y$};\draw[blue,thick,solid,domain=0.03:2.15,samples=180] plot(\x,{(\x)^2});\draw[teal,thick,dashed,domain=0.03:2.6,samples=180] plot(\x,{\x});\draw[violet,thick,densely dotted,domain=0.03:2.6,samples=180] plot(\x,{sqrt(\x)});\fill[orange!85!black] (1,1) circle(2.2pt);\draw[gray,dashed] (1,0)--(1,4.6);\node[below=5pt] at (1,0){$1$};\end{scope}\node[anchor=west,align=left,text width=5.0cm] at (8.4,4.0) {蓝实线：$x^2$\\[8pt]青虚线：$x$\\[8pt]紫点线：$\sqrt{x}$\\[8pt]$0<x<1$：$x^2<x<\sqrt{x}$\\[8pt]$x>1$：$x^2>x>\sqrt{x}$};
\node[text width=12.6cm,align=center] at (7.1,.65) {这里始终固定系数为1；三条曲线在正半轴上只有公共点(1,1)。};
\end{tikzpicture}\end{center}



  把“提高指数”想成再乘上一个相同的数：
  \[
     2,\ 2^2=4,\ 2^3=8;\qquad
     \tfrac12,\ (\tfrac12)^2=\tfrac14,\ (\tfrac12)^3=\tfrac18.
  \]
  就像每轮翻倍的积分与每轮只剩一半的钱，前者轮数越多越大，后者轮数越多越小。整数指数只是帮助建立直觉，分数、负数和其他实数指数也满足同样的比较规律（见证明）。
  看虚线 $x=1$：右边指数大的在上方，左边恰好反过来。这里只画 $x>0$ 的部分；是否能取 $x=0$，留到易错三单独看。

  \medskip
  
\end{explanationbox}
\clearpage
\begin{explanationbox}
\textbf{\textcolor{CExplanation}{六、一个特别有用的发现：在公共点处，指数就是切线斜率}}
\begin{center}\begin{tikzpicture}[font=\small,>=stealth,line join=round]
\draw[gray!40,rounded corners=3pt] (0,0) rectangle (14.2,8.5);\node at (7.1,7.9) {在公共点处，指数直接给出切线的斜率};
\begin{scope}[shift={(0.7,2.6)},x=1.5cm,y=1.05cm]\draw[->,gray] (-.15,0)--(2.25,0);\draw[->,gray] (0,-.3)--(0,3.3);\draw[blue,thick,domain=.06:1.75,samples=120] plot(\x,{(\x)^(2)});\draw[orange!85!black,dashed,domain=.45:1.55] plot(\x,{1+(2)*(\x-1)});\fill[orange!85!black] (1,1) circle(2pt);\end{scope}\node at (2.4,6.9){$x^2$};\node at (2.4,1.7){$f^{\prime}(1)=2$};\begin{scope}[shift={(5.3500000000000005,2.6)},x=1.5cm,y=1.05cm]\draw[->,gray] (-.15,0)--(2.25,0);\draw[->,gray] (0,-.3)--(0,3.3);\draw[teal,thick,domain=.06:1.75,samples=120] plot(\x,{(\x)^(1)});\draw[orange!85!black,dashed,domain=.45:1.55] plot(\x,{1+(1)*(\x-1)});\fill[orange!85!black] (1,1) circle(2pt);\end{scope}\node at (7.050000000000001,6.9){$x$};\node at (7.050000000000001,1.7){$f^{\prime}(1)=1$};\begin{scope}[shift={(10.0,2.6)},x=1.5cm,y=1.05cm]\draw[->,gray] (-.15,0)--(2.25,0);\draw[->,gray] (0,-.3)--(0,3.3);\draw[violet,thick,domain=.06:1.75,samples=120] plot(\x,{(\x)^(0.5)});\draw[orange!85!black,dashed,domain=.45:1.55] plot(\x,{1+(0.5)*(\x-1)});\fill[orange!85!black] (1,1) circle(2pt);\end{scope}\node at (11.700000000000001,6.9){$\sqrt{x}$};\node at (11.700000000000001,1.7){$f^{\prime}(1)=0.5$};
\node[text width=12.6cm,align=center] at (7.1,.65) {三图同尺度；橙虚线为(1,1)处真实切线，斜率分别为2、1、二分之一。};
\end{tikzpicture}\end{center}



  所有基本幂函数都经过 $(1,1)$。学过导数后会发现，$f'(1)=\alpha$。例如，$x^2$ 在此处的切线斜率为 $2$，$x$ 为 $1$，$\sqrt{x}$ 为 $1/2$，$1/x$ 为 $-1$。所以不仅能从函数式判断图像走势，也能根据图像在公共点处的切线斜率反推出指数。这是系数固定为 $1$ 带来的另一个好处。

  \medskip
  
\end{explanationbox}
\clearpage
\begin{explanationbox}
\textbf{\textcolor{CExplanation}{七、一个容易忽略的联系：幂函数与指数函数}}
\begin{center}\begin{tikzpicture}[font=\small,>=stealth,line join=round]
\draw[gray!40,rounded corners=3pt] (0,0) rectangle (14.2,8.8);\node at (7.1,8.200000000000001) {变量放在不同位置，就是不同的研究对象};
\node at (3.1500000000000004,7.3) {$y=x^2$};\begin{scope}[shift={(1.2444444444444445,2.7)},x=1.3611111111111112cm,y=0.39999999999999997cm]\draw[->,gray] (-0.4,0)--(3.2,0) node[right]{$x$};\draw[->,gray] (0,-0.5)--(0,9) node[above]{$y$};\draw[blue,thick,solid,domain=0:2.8,samples=180] plot(\x,{(\x)^2});\fill[orange!85!black] (1,1) circle(2pt);\fill[orange!85!black] (2,4) circle(2pt);\end{scope}\node[align=center,text width=5.8cm] at (3.1500000000000004,1.7) {指数固定为2\\改变底数x};\node at (9.95,7.3) {$y=2^x$};\begin{scope}[shift={(8.044444444444444,2.7)},x=1.3611111111111112cm,y=0.39999999999999997cm]\draw[->,gray] (-0.4,0)--(3.2,0) node[right]{$x$};\draw[->,gray] (0,-0.5)--(0,9) node[above]{$y$};\draw[orange!85!black,thick,solid,domain=-0.2:2.8,samples=180] plot(\x,{2^\x});\fill[orange!85!black] (1,2) circle(2pt);\fill[orange!85!black] (2,4) circle(2pt);\end{scope}\node[align=center,text width=5.8cm] at (9.95,1.7) {底数固定为2\\改变指数x};
\node[text width=12.6cm,align=center] at (7.1,.65) {两图同尺度；同底幂比较可以转成指数函数的单调性问题。};
\end{tikzpicture}\end{center}



  同一个式子 $x^\alpha$，固定指数、改变底数，就是幂函数；固定底数、改变指数，就是我们学过的指数函数（底数为 $1$ 时是常数）。所以比较同底幂时，可以借用指数函数的增减规律，而不必算出具体数值。

  \medskip
  
\end{explanationbox}
\clearpage
\begin{explanationbox}
\textbf{\textcolor{CExplanation}{八、再进一步：为什么负指数图像靠近坐标轴？}}
\begin{center}\begin{tikzpicture}[font=\small,>=stealth,line join=round]
\draw[gray!40,rounded corners=3pt] (0,0) rectangle (14.2,8.3);\node at (7.1,7.700000000000001) {负指数：越靠近轴，不等于真的碰到轴};
\begin{scope}[shift={(1.6,2.033333333333333)},x=1.0cm,y=0.7666666666666666cm]\draw[->,gray] (-0.5,0)--(5.5,0) node[right]{$x$};\draw[->,gray] (0,-0.5)--(0,5.5) node[above]{$y$};\draw[orange!85!black,thick,solid,domain=0.2:5,samples=180] plot(\x,{1/\x});\fill[orange!85!black] (0.25,4) circle(2.2pt);\fill[orange!85!black] (4,0.25) circle(2.2pt);\node[right=5pt] at (.25,4){$(\frac14,4)$};\node[above=5pt] at (4,.25){$(4,\frac14)$};\end{scope}\node[anchor=west,align=left,text width=5.0cm] at (8.4,4.0) {$y=1/x,\ x>0$\\[9pt]$(\frac14,4)$与$(4,\frac14)$\\[9pt]$x\to0^+$：$y\to+\infty$\\[9pt]$x\to+\infty$：$y\to0^+$\\[9pt]始终有$x>0,y>0$};
\node[text width=12.6cm,align=center] at (7.1,.65) {图取负指数中的1/x；一般x的负p次方具有相同的正半轴渐近行为，p为正数。};
\end{tikzpicture}\end{center}



  写成 $x^{-p}=1/x^p$（$p>0$）。当 $x$ 越来越大，分母越来越大，结果就越来越接近 $0$；当 $x$ 从右边越来越靠近 $0$，分母越来越小，结果就越来越大。例如 $1/100=0.01$，$1/0.01=100$。但只要 $x>0$，函数值始终为正，不会碰到坐标轴。
\end{explanationbox}
